
\section{Removing Excess Items by Greedy Deflation}
\label{app:greedy-deflation}

We describe a greedy deflation procedure for
converting an integer solution satisfying the
covering constraints into a packing containing
exactly the required items.

Let $\mathbf{b}=\hat{\mathbf{a}}$ denote the
configuration multiplicities after rounding.
For each configuration
$C_i=(c_{i1},\ldots,c_{im})$, define
$C_i^{(j\leftarrow x)}$ as the configuration
obtained by replacing only its $j$-th
coordinate with $x$.

\begin{algorithm}[H]
    \caption{Greedy Deflation}
    \label{alg:greedy-deflation}
    \KwIn{Feasible configurations $C_1,\ldots,C_q$, required item
        counts $n_1,\ldots,n_m$, and integer multiplicities
        $\hat{\mathbf{a}}$ satisfying the covering constraints}
    \KwOut{Configuration multiplicities $\mathbf{b}$ packing exactly
        the required items, with no increase in the number of bins}
    $\mathbf{b} \gets \hat{\mathbf{a}}$\;
    \For{$j=1,\ldots,m$}{
        $e_j \gets \sum_{i=1}^{q}b_i c_{ij}-n_j$\;
        \While{$e_j>0$}{
            Choose a configuration $C_i$ with $b_i>0$ and $c_{ij}>0$\;
            $t \gets c_{ij}$, $d \gets \min\{e_j,b_i t\}$\;
            $k \gets \lfloor d/t\rfloor$, $r \gets d\bmod t$\;
            Replace $k$ bins of configuration $C_i$
                with $C_i^{(j\leftarrow 0)}$\;
            \If{$r>0$}{
                Replace one additional bin of configuration $C_i$
                    with $C_i^{(j\leftarrow t-r)}$\;
            }
            Update the corresponding multiplicities in $\mathbf{b}$\;
            $e_j \gets e_j-d$\;
            Discard any empty bins\;
        }
    }
    \Return{$\mathbf{b}$}\;
\end{algorithm}

The algorithm processes excess items in batches
rather than examining individual bins.


\paragraph{Illustrative Example.}

Suppose the item sizes are
$y_1=0.1$ and $y_2=0.2$.
The rounded solution contains 5 bins of
configuration $(3,2)$, whereas the instance
requires $n_1=8$ and $n_2=7$ items.

Initially, the selected configurations contain
\[
    5(3,2)=(15,10).
\]

For the first item type, the excess is
$e_1=15-8=7$. We obtain
\[
    d=7,\qquad k=2,\qquad r=1.
\]

Hence, two bins are deflated to $(0,2)$
and one bin to $(2,2)$:
\[
    5(3,2)
    \longrightarrow
    2(3,2)+2(0,2)+(2,2).
\]

The resulting configurations contain
exactly eight items of the first type,
but ten items of the second type.

For the second item type, $e_2=3$.
Selecting the two bins of configuration
$(0,2)$ gives
\[
    d=3,\qquad k=1,\qquad r=1.
\]

One bin becomes empty and is discarded,
while the other is deflated to $(0,1)$:
\[
    2(3,2)+2(0,2)+(2,2)
    \longrightarrow
    2(3,2)+(2,2)+(0,1).
\]

The resulting packing contains exactly
\[
    2(3,2)+(2,2)+(0,1)=(8,7),
\]
using only four bins.

% Insert a TikZ illustration of the
% configuration transformations here.


\paragraph{Correctness.}
For each item type $j$, the covering constraints
guarantee that at least $n_j$ items are available.
Thus, the algorithm can remove exactly $e_j$
excess items.

By definition, the deflation transformation
$C_i^{(j\leftarrow x)}$ modifies only the
$j$-th coordinate of configuration $C_i$,
leaving all other coordinates unchanged.
Consequently, for every item type $\ell\neq j$,
the total number of type-$\ell$ items remains
unchanged. In particular, the equality constraints
established for previously processed item types
are preserved.

Furthermore, since deflation only removes items
from existing bins, it preserves bin feasibility
without increasing the number of bins.



\paragraph{Running Time.}
For each of the $m$ item types, the algorithm
processes at most $q$ configurations.
Each configuration is processed using a
constant number of arithmetic operations,
assuming configuration transitions can be
located in constant time.

Therefore, deflation requires
\[
    O(mq)
\]
arithmetic operations, excluding the
preprocessing required to index configurations.

